We provide here a theoretical foundation for the objective in~\Eqref{eq:mle ultimate form}, showing how it arises naturally from a maximum likelihood estimation (MLE) framework under mild assumptions.

Let $\vx \in \mathbb{R}^d$ be a random vector (considered as a token in the context of Transformer) in a high-dimensional representation space with a probability distribution $p(\vx)$. Let $\{\vz^h\}_{h=1}^H$ be a set of random vectors in low-dimensional latent spaces $\mathbb{R}^{p} ~(p<d)$ with distinct support, following a prior joint probability distribution $p(z^1, \dots, \vz^H)$. We formulate information processing in the forward pass of neural networks as maximum likelihood estimation:
\begin{equation}
\label{eq:mle}
    \max_{\vx} \quad \mathbb{E}_{(\vz^1, \dots,\vz^H)\sim p(\vz^1,\dots,\vz^H)}\left[\log p(\vx,\vz^1,\dots,\vz^H; \theta,\phi)\right],
\end{equation}
where $\theta$ and $\phi$ are parameters of the high- and low-dimensional encodings, respectively.

To make this optimization more tractable, we make the following basic and practical assumptions:
\begin{assumption}
\label{assumption:1}
The random vectors $\vz^1,\dots,\vz^H$ are independent and follow the identical distribution $p(\vz)$ in distinct latent spaces, \ie, $p(\vz^1,\dots,\vz^H)=\Pi_{h=1}^H p(\vz^h)$ and $p(\vz^1)=\dots=p(\vz^H)=p(\vz)$.
\end{assumption}

\begin{assumption}
\label{assumption:2}
The prior distribution $p(\vz)$ is a uniform distribution with support on a hypersphere $\mathbb{S}^{p-1}$.
\end{assumption}

\begin{assumption}
\label{assumption:3}
The random vectors $(\vz^1,\dots,\vz^H)\sim p_{\phi}(\vz^1,\dots,\vz^H \mid \vx)$ from the posterior distribution are conditionally independent, \ie, $p_{\phi}(\vz^1,\dots,\vz^H \mid \vx) = \Pi_{h=1}^H p_\phi(\vz^h \mid \vx)$.
\end{assumption}

Assumption~\ref{assumption:2} of hyperspherical uniform distribution can be perceived to function as regularization on the latent representations to preserve maximum entropy and avoid representational collapse, which has been adopted to enhance auto-encoding \citep{xu-durrett-2018-spherical,s-vae18}. Under the above basic and practical assumptions, the MLE objective can be reformulated as:
\begin{align}
\label{eq:mle reformulation}
    \max_{\vx}~& \mathbb{E}_{(\vz^1, \dots,\vz^H)\sim p(\vz^1,\dots,\vz^H)}\left[\log p(\vx,\vz^1,\dots,\vz^H; \theta,\phi)\right] \nonumber \\
    & = \mathbb{E}_{(\vz^1,\dots,\vz^H) \sim p(\vz)}\left[ \log p_\phi(\vz^1,\dots,\vz^H \mid \vx)\right] + \mathbb{E}_{(\vz^1,\dots,\vz^H) \sim p(\vz)} \left[\log p_\theta(\vx) \right]  \nonumber\\
    & = \sum_{h=1}^H \mathbb{E}_{\vz^h \sim p(\vz)} \left[\log p_\phi(\vz^h \mid \vx) \right] + \log p_\theta(\vx) \nonumber \\
    & = \sum_{h=1}^H \mathbb{E}_{\vz^h \sim p(\vz)} \left[\log \frac{p_\phi(\vz^h \mid \vx)}{p(\vz^h)} \right] + \sum_{h=1}^H \mathbb{E}_{\vz^h \sim p(\vz)} \left[\log  p(\vz^h) \right] + \log p_\theta(\vx) \nonumber \\
    & = \sum_{h=1}^H \mathbb{E}_{\vz^h \sim p(\vz)} \left[\log \frac{p_\phi(\vz^h \mid \vx)}{p(\vz)} \right] + \sum_{h=1}^H \mathbb{E}_{\vz \sim p(\vz)} \left[\log p(\vz) \right] + \log p_\theta(\vx) \nonumber \\
    & = -\sum_{h=1}^H \KL( p(\vz) \Vert p_\phi(\vz^h \mid \vx) ) - H \times \mathcal{H}(p(\vz)) +  \log p_\theta(\vx), 
\end{align}
where $\KL(\cdot \Vert \cdot)$ denotes Kullback-Leibler (KL) divergence and $\mathcal{H}(\cdot)$ means differential entropy. As the second term on entropy in~\Eqref{eq:mle reformulation} does not depend on variable $\vx$, this objective ultimately reduces to~\Eqref{eq:mle ultimate form} which we restate here for completeness:
\begin{equation*}
% \label{eq:mle ultimate form}
\min_{\vx} \quad \sum_{h=1}^H \underbrace{\KL \left( p(\vz) \Vert p_\phi(\vz^h \mid \vx) \right)}_{\text{uniformity}} \underbrace{- \log(p_\theta(\vx))}_{\text{alignment}}.
\end{equation*}

The first term encourages the posterior $p_\phi(\vz^h \mid \vx)$ defined on the vector $\vz^h \in \mathbb{R}^p$ in latent space, which can be implemented by a transformation parameterized by $\phi$, to approximate a uniform distribution on a hypersphere. To see why the second term implies alignment, suppose the distribution $p_\theta(\vx)$ is parameterized as a mixture of $M$ von Mises–Fisher (vMF) distributions\footnote{\href{https://en.wikipedia.org/wiki/Von_Mises–Fisher_distribution}{\text{https://en.wikipedia.org/wiki/Von\_Mises–Fisher\_distribution}}} with equal mixing coefficients:
\begin{equation}
\label{eq:vMF}
-\log(p_\theta(\vx)) = -\log\left(\frac{1}{M} \sum_{m=1}^M f(\vx; \vmu_m, \kappa_m) \right)= -\log\left(\frac{1}{M} \sum_{m=1}^M C_{d}(\kappa_m)\exp(\kappa_m \vmu_m^\top\vx)\right),
\end{equation}
where $\vmu_m \in \mathbb{S}^{d-1}$ denotes mean direction on $(d-1)$-dimensional unit sphere and $\kappa_m \geq 0$ is the concentration parameter while $C_{d}(\kappa_m)$ is the normalization constant; therefore, finding the vector $\vx$ that minimizes this negative log-probability~\Eqref{eq:vMF} equals finding maximal inner product $\vmu_m^\top\vx$, thus aiming for directional alignment. In practice, the mean direction $\vmu_m$ is learned by backpropagation and consequently contains certain statistical properties from the data. 

In summary, the objective~\Eqref{eq:mle ultimate form} suggests that for a representation vector $\vx \in \mathbb{R}^d$, the forward dynamics can be characterized by two complementary properties: 
\begin{itemize}[leftmargin=*]
    \item \textbf{Mode Seeking}: Achieving semantic alignment with directional vectors encapsulating specific information derived from data in the \textbf{high-dimensional space}.
    \item \textbf{Mass Covering}: Maximally preserving the entropy embedded via regularizing distributional uniformity in the \textbf{low-dimensional space}.
\end{itemize}
These principles underpin our design of token dynamics, and we propose to use energy functions to quantify these two properties as instantiations that can induce various Transformer-based models. 